% !TeX program = pdflatex
% !TeX encoding = UTF-8
% 根据 docs/template/main.tex 的版式与章节结构创建。
% 正式竞标书；中英文内容同步维护。
% 编译：使用 pdfLaTeX，连续编译两次更新目录。
% 成员资料来自 docs/character/；不得用往年成员信息填补未知字段。
\documentclass[a4paper,12pt]{article}
\usepackage[utf8]{inputenc}
\usepackage{amsmath}
\usepackage{array}
\usepackage[table]{xcolor}
\usepackage{needspace}
\usepackage{fancyhdr}
\usepackage{graphicx}
\usepackage[hidelinks]{hyperref}
\usepackage{tocloft}
\renewcommand{\labelitemi}{\ensuremath{\bullet}}
\setlength{\headheight}{14.5pt}
\addtolength{\headsep}{-2.5pt}
% Allow a little extra spacing on difficult lines without changing template margins.
\setlength{\emergencystretch}{3em}

% 统一使用本组名称，不沿用往年模板的组号。
\newcommand{\BidTeam}{UNNC GRP Team202610}
\newcommand{\BidDate}{October 15, 2026}
\newcommand{\BidHeaderDate}{October 15, 2026}
\newcommand{\PersonalInformation}[2]{%
  \subsection*{Personal Information}
  \if\relax\detokenize{#1}\relax\else
  \noindent\textbf{Student ID:} #1\\
  \fi
  \textbf{Email:} #2\par
}
\pagestyle{fancy}
\fancyhf{}
\fancyhead[L]{\BidTeam}
\fancyhead[R]{\BidHeaderDate}
\fancyfoot[C]{\thepage}
\renewcommand{\headrulewidth}{0.4pt}
\renewcommand{\footrulewidth}{0pt}

\newcommand{\BidProjectCode}{P19}
\newcommand{\BidProjectName}{SmartCareer: A Personalized Employment Recommendation System with a Multi-Dimensional Job Quality Evaluation Model}
\newcommand{\BidTitle}{Bid for [\BidProjectCode]}
\title{\BidTitle}
\author{\BidTeam}
\date{\BidDate}
\hypersetup{pdftitle={Bid for [P19]},pdfauthor={UNNC GRP Team202610}}

\begin{document}
\maketitle
\thispagestyle{plain}
\tableofcontents
\newpage

\section{Introduction}
\noindent\textbf{Project:} \BidProjectName
\par\medskip
UNNC GRP Team202610 brings together experience in data preparation, model evaluation, and interactive software development that fits SmartCareer's main challenges. Yun Xin NG's data-engineering internship and machine-learning projects support working with recruitment data, while Ziqian XIONG's model experiments and Flask/React projects provide a basis for connecting recommendation methods with a usable application. Research annotation, programming, database, and CI/CD skills across the team complement these experiences.

As university students, we want to help our peers understand why a job suits them and which skills they could develop. We will combine student profiles and skill matching with transparent scores for salary, growth, stability, and skill fit, so that users can explore how their priorities change the recommendations.

Led by Ziqian XIONG, we will first demonstrate an explainable recommendation workflow on the Zhejiang graduate subset, then evaluate trend forecasting, occupation relationships, and ranking improvements. Shared learning, regular integration, and early supervisor feedback will help us turn our complementary preparation into a working system. We will assess recommendations through baseline comparisons and student trials, and clearly show the dates and limitations of historical job data.

\newpage
\section{Teammates}
Our seven-member team, led by Ziqian XIONG, combines data preparation and model evaluation, algorithm implementation, application and database development, and testing, deployment, and user communication. We will connect these strengths through shared data definitions, module interfaces, and integration reviews so that experimental results become demonstrable software. The following profiles present the experience supporting our project plan; responsibilities are set out in the work-distribution section.

We believe that the bidding groups have broadly comparable academic foundations. We therefore see stronger leadership and proactive engagement as our competitive advantage: a leader who can guide shared learning and coordinate decisions, and members willing to investigate unfamiliar methods, raise difficulties early, and follow through on feedback. We will put this into practice through weekly progress reviews during concentrated development, paired method discussions, and early demonstrations; the low-workload period will use brief asynchronous check-ins.

\subsection{Ziqian XIONG}
\PersonalInformation{20721428}{\href{mailto:scyzx7@nottingham.edu.cn}{\nolinkurl{scyzx7@nottingham.edu.cn}}}
\noindent\textbf{Group role:} Leader\par

\subsection*{Introduction}
Ziqian XIONG is a Computer Science undergraduate at UNNC and our team leader. His medical image segmentation research involves model implementation, dataset checks, and comparative evaluation. His 4DGS tool and MomoFocus application provide Python/Flask and React/TypeScript development experience. For P19, these experiences provide a basis for comparing skill-matching and ranking methods and connecting them to an interactive recommendation application. Personal website: \url{https://qqqqqprofile.vercel.app/}.

\subsection*{Internship Experience}
\begin{itemize}
  \item \textbf{Pony.ai algorithm internship (July--September 2026):} In the AI department, investigated VGGT-related 3D/4D reconstruction methods and compared multi-view geometry and depth estimation against an existing DA3/3D Gaussian Splatting pipeline. Completed feasibility validation and engineering integration, replacing the original reconstruction workflow. Also wrote documentation on agent-assisted development workflows.
\end{itemize}

\subsection*{Project Experience}
\begin{itemize}
  \item \textbf{4DGS-Edit-and-Compare (July--September 2026):} Independently developed a 4D Gaussian Splatting point-cloud editing and comparison tool. It supports rotation, translation, and deletion of selected regions, editing and comparing multiple point clouds, and quantitative evaluation. \textbf{Technology stack:} Python/Flask for the backend, JavaScript/Three.js/WebGL for the interactive frontend, NumPy/SciPy for numerical processing, PyTorch/plyfile for point-cloud file handling, and Docker for packaging. Developed independently during the Pony.ai internship, the tool is reported on his website as being used by the company. \href{https://github.com/qianqqqqqXZQ/4DGS-Edit-and-Compare}{Project repository}.
  \item \textbf{MomoFocus (September--October 2026):} Developed a personal productivity application integrating to-do lists, notes, and a Pomodoro timer. \textbf{Technology stack:} React, TypeScript, and CSS for the interface; Electron for the desktop shell; Vite for development and frontend builds; and electron-builder for desktop packaging. The implementation uses localStorage for local persistence and dnd-kit for sortable lists, providing relevant experience in interactive application development and desktop delivery. \href{https://github.com/qianqqqqqXZQ/MomoFocus}{Project repository}.
\end{itemize}

\subsection*{Research Experience}
\begin{itemize}
  \item \textbf{Medical image segmentation research:} Implemented and debugged CNN--Transformer components; prepared and checked datasets; evaluated Dice, IoU, and inference speed. \textbf{Technology stack:} Python, PyTorch, and OpenMMLab's MMSegmentation, MMCV, and MMEngine. A related paper is under submission to CVIU with him as second author.
  \item \textbf{Dynamic dual-branch neural network research:} Investigating a multi-task network for keyword spotting and speaker verification, with early-exit scheduling and FPGA resource reallocation. Network design, scheduling, and resource-allocation strategies remain under evaluation.
\end{itemize}

\subsection*{Skill Set}
\begin{itemize}
  \item \textbf{Programming languages:} Python, C/C++, Java, Haskell, JavaScript, TypeScript, Lean, and Verilog.
  \item \textbf{Databases:} MySQL.
  \item \textbf{Research and implementation:} PyTorch; OpenMMLab (MMSegmentation, MMCV, and MMEngine); CNN--Transformer models; data preprocessing and comparative model evaluation; Retrieval-Augmented Generation (RAG).
  \item \textbf{Application development:} HTML, CSS, React, Electron, Vite, Flask, Three.js/WebGL, NumPy, and SciPy.
  \item \textbf{Engineering and documentation:} Docker packaging, electron-builder, and technical documentation for agent-assisted development workflows, including harness and loop engineering.
\end{itemize}

\subsection{Yun Xin NG}
\PersonalInformation{20709091}{efyyn8@nottingham.edu.cn}
\noindent\textbf{Group role:} Member\par

\subsection*{Introduction}
Yun Xin NG is a second-year Computer Science student on a one-year exchange at UNNC after completing her first year at the University of Nottingham Malaysia. Her YTL AI Labs internship involved multi-source data extraction and benchmark-based model selection, while her machine learning projects involved dataset splitting and training. For P19, this experience provides a basis for cleaning recruitment data, preparing evaluation samples, and comparing skill-trend forecasts.

\subsection*{Internship Experience}
\begin{itemize}
  \item \textbf{YTL AI Labs data-engineering internship:} Extracted data from web pages and PDF documents, shortlisted open-source vision-language models against application needs, and compared document-understanding and text-extraction capabilities using existing benchmarks. Selected a model by considering both benchmark performance and practical suitability.
\end{itemize}

\subsection*{Project Experience}
\begin{itemize}
  \item \textbf{Machine learning project:} Prepared training/test splits and trained models, gaining experience in dataset preparation and the experimental workflow.
  \item \textbf{Internship Management System:} Developed a first-year course application using HTML, CSS, Java, and MySQL, connecting page development with application logic and persistent data.
\end{itemize}

\subsection*{Skill Set}
\begin{itemize}
  \item \textbf{Programming languages:} Python, Java, and C/C++.
  \item \textbf{Web and database technologies:} HTML, CSS, and MySQL.
  \item \textbf{Data engineering:} Data extraction from web pages and PDFs, dataset preparation, training/test splitting, and Retrieval-Augmented Generation (RAG).
  \item \textbf{Models and AI technologies:} Model training and benchmark-based vision-language model evaluation and selection.
\end{itemize}

\subsection{Jinyao LU}
\PersonalInformation{}{scyl@nottingham.edu.cn}
\noindent\textbf{Group role:} Member\par

\subsection*{Introduction}
Jinyao LU is a Computer Science and Artificial Intelligence undergraduate at UNNC with LLM-related research, technical-writing, and software-workflow experience. His internships involved platform business logic, information-platform maintenance, and data recording; he also helped organise an AI hackathon. For P19, these experiences support organising student requirements and trials, checking recommendation explanations, and preparing documentation and demonstrations.

\subsection*{Internship Experience}
\begin{itemize}
  \item \textbf{R\&D internship at Zhejiang Dishi Information Technology (July--August 2024):} Learned the software-development workflow and platform business logic; investigated VR/AR applications in the water sector.
  \item \textbf{Internship at Zhejiang Zhongkong Information Industry (June--August 2025):} Monitored train status, participated in routine maintenance of an information-publishing platform and operational data recording, and learned to operate the dispatch system.
\end{itemize}

\subsection*{Project Experience}
\begin{itemize}
  \item \textbf{Douyin AI hackathon organisation (March--April 2026):} Served on the campus core organising team, coordinating participating teams and event processes.
\end{itemize}

\subsection*{Research Experience}
\begin{itemize}
  \item \textbf{AI research:} First author of a paper on LLMs and autonomous driving accepted by an international journal.
\end{itemize}

\subsection*{Skill Set}
\begin{itemize}
  \item \textbf{Programming languages:} Python, Java, and C/C++.
  \item \textbf{Databases:} MySQL.
  \item \textbf{Technologies and foundations:} NumPy data processing, basic web development, data structures, algorithms, and software engineering.
  \item \textbf{Development and reporting tools:} Environment setup and basic use of C and Java programming tools; Excel charts and PowerPoint presentations.
  \item \textbf{Technical communication:} Academic writing and presentation of research findings.
\end{itemize}

\subsection{Oswald Yiwen CHONG}
\PersonalInformation{20792699}{\href{mailto:hcyoc1@nottingham.edu.cn}{\nolinkurl{hcyoc1@nottingham.edu.cn}}}
\noindent\textbf{Group role:} Member\par

\subsection*{Introduction}
Oswald Yiwen CHONG is from Malaysia and studies Computer Science with Artificial Intelligence at the University of Nottingham. His blockchain-based bouldering social application provides practical application-development experience, complemented by GitHub collaboration and CI/CD skills. For P19, these capabilities support connecting student profiles, recommendation services, and feedback into a tested application with repeatable deployment.

\subsection*{Project Experience}
\begin{itemize}
  \item \textbf{Blockchain-based bouldering social application:} Developed a Strava-like social application for bouldering, using blockchain to award user badges and incorporating NFT-related functionality. This provided practical experience in applying decentralised technology to a user-facing application.
\end{itemize}

\subsection*{Skill Set}
\begin{itemize}
  \item \textbf{Programming languages:} Python, Java, and C/C++.
  \item \textbf{Databases:} MySQL.
  \item \textbf{Application technologies:} Application development, blockchain, and decentralised technologies.
  \item \textbf{Version control and collaboration:} Proficient use of GitHub for version control and team collaboration.
  \item \textbf{Continuous integration and continuous deployment (CI/CD):} Able to set up pipelines that support code reliability and rapid software delivery.
\end{itemize}

\subsection{Ziqian ZHANG}
\PersonalInformation{20721592}{scyzz23@nottingham.edu.cn}
\noindent\textbf{Group role:} Member\par

\subsection*{Introduction}
Ziqian ZHANG is a third-year Computer Science student with research experience in image annotation and data preparation, alongside Labelme and JSON data-management skills. For P19, we will adapt this experience to preparing checked skill and occupation samples, organising annotations, and recording recruitment-data and recommendation errors.

\subsection*{Research Experience}
\begin{itemize}
  \item \textbf{AI-driven Analysis and Optimization of Performance on Magnets:} Participated in this research project, contributing image annotation and data preparation. Alongside this experience, her skills in Labelme and JSON annotation/data management provide a foundation for organising image labels and structured annotation records. These capabilities can support the preparation of labelled examples and evaluation materials, with attention to consistency between source images and their annotations.
\end{itemize}

\subsection*{Skill Set}
\begin{itemize}
  \item \textbf{Programming languages:} Python, Java, and C/C++ at a foundational level.
  \item \textbf{Databases:} MySQL.
  \item \textbf{Data tools and processing:} Labelme, JSON annotation and data management, image annotation, and dataset preparation.
\end{itemize}

\subsection{Yue WU}
\PersonalInformation{20717823}{\href{mailto:sgyyw10@nottingham.edu.cn}{\nolinkurl{sgyyw10@nottingham.edu.cn}}}
\noindent\textbf{Group role:} Member\par

\subsection*{Introduction}
Yue WU studies Computer Science at UNNC. Her programming coursework and CCN research writing provide experience in implementation, literature review, and technical analysis. For P19, her Python and database foundations support developing skill-normalisation rules, four-dimensional scoring, and occupation relationships, with documented assumptions and method comparisons.

\subsection*{Project Experience}
\begin{itemize}
  \item \textbf{Coursework development:} Applied programming and algorithmic problem-solving in software-related course projects, building experience in translating computing concepts into implementations.
\end{itemize}

\subsection*{Research Experience}
\begin{itemize}
  \item \textbf{CCN research and technical writing:} Worked on research and paper writing concerning \emph{The softwarization and interoperability of CCN}, developing literature-review, technical-analysis, and academic-writing experience.
\end{itemize}

\subsection*{Skill Set}
\begin{itemize}
  \item \textbf{Programming languages:} Java, Python, Haskell, C/C++, and JavaScript.
  \item \textbf{Web and database technologies:} HTML, CSS, MySQL, and SQLite.
  \item \textbf{Development and analysis tools:} IntelliJ IDEA, Anaconda, Visual Paradigm, LaTeX, and MATLAB.
  \item \textbf{Technical analysis:} Algorithm analysis, technical literature review, and academic writing.
\end{itemize}

\subsection{Xinyun HU}
\PersonalInformation{20720227}{scyxh8@nottingham.edu.cn}
\noindent\textbf{Group role:} Member\par

\subsection*{Introduction}
Xinyun HU is a third-year Computer Science student with programming, SQL, database, and web-technology foundations, alongside an AI-enabled Connect Four project. For P19, these skills provide a basis for building student-profile and preference forms, presenting ranked jobs and skill gaps, and checking filtering, weight changes, and feedback interactions.

\subsection*{Project Experience}
\begin{itemize}
  \item \textbf{Connect Four Game with AI (July 2025):} The supplied project background includes an AI-enabled Connect Four game, providing a relevant setting for discussing the interaction between algorithmic behaviour and an interactive application.
\end{itemize}

\subsection*{Skill Set}
\begin{itemize}
  \item \textbf{Programming and query languages:} Java, Python, C/C++, Haskell, SQL, MATLAB, and MIPS Assembly.
  \item \textbf{Web and database technologies:} HTML, CSS, MySQL, and SQLite.
  \item \textbf{Development tools:} Git, VS Code, Eclipse, and Visual Paradigm.
\end{itemize}

\subsection{Why Our Team Fits This Project}
\begin{itemize}
  \item \textbf{Data quality before recommendation complexity:} Yun Xin NG's dataset-preparation and model-training experience and Ziqian ZHANG's research annotation background support a documented recruitment-data pipeline and checked evaluation samples. This addresses the need to make real job descriptions, skills, and salary fields usable before ranking.
  \item \textbf{From a transparent baseline to a usable service:} Ziqian XIONG's comparative model evaluation, Yue WU's Python/SQLite and analysis background, and Xinyun HU's programming and database foundations provide a basis for developing and checking skill matching, quality scoring, and profile-to-result workflows. Oswald Yiwen CHONG can support GitHub collaboration, automated checks, and deployment; Jinyao LU can support student-facing requirements, trials, and demonstrations.
  \item \textbf{A focused seven-person delivery structure:} We will organise data, algorithms, and application workstreams with explicit handovers and shared integration checkpoints. We will first demonstrate a transparent baseline on the Zhejiang graduate subset, then evaluate trends, graph features, and ranking improvements. Our seven members will share implementation and review across these three workstreams, with one leader coordinating interfaces and delivery.
\end{itemize}

\newpage

\section{Ideas \& Design}
\subsection{Project Proposal}
SmartCareer will help students understand which occupations fit their skills, why a job is recommended, and which skills they need to develop. Students will enter their major, skills, qualifications, location, and career priorities, adjust the importance of salary, growth, stability, and skill fit, and receive ranked jobs with dimension scores, explanations, and skill gaps. An occupation--skill graph will show related career pathways, while explicit relevance feedback will support user-approved preference updates.

Development will begin with the Zhejiang graduate subset of the research group's recruitment data and scale through batch processing. Historical records will support labour-market analysis and recommendation evaluation; the interface will display their dates and distinguish them from current vacancies. The core deliverables are a clean queryable dataset, skill dictionary and new-occupation graph, documented quality model, evaluated recommendation module, working web application, and final report, source code, and demonstration.

\subsection{Technology Stack \& System Design}
\begin{enumerate}
  \item \textbf{Back-end Development}\par
  Python and pandas will clean records in batches, standardise salary periods and ranges, remove cross-platform duplicates, and preserve source and transformation records. Parquet will hold offline features, DuckDB will query recruitment data, and SQLite will store student profiles, preferences, and feedback. FastAPI will expose profile, search, recommendation, explanation, graph, and feedback endpoints.

  Skill dictionaries and synonym rules will provide an extraction baseline. scikit-learn will support keyword and skill matching, hybrid ranking, and forecasting comparisons. Monthly or quarterly skill shares will feed seasonal-naive and lag-feature regression forecasts, evaluated using chronological splits. NetworkX will represent occupation--skill relationships; persistent changes in titles, skill combinations, and demand will identify candidate emerging occupations for manual review.

  The four quality dimensions will use comparable salary percentiles, demand and skill trends, demand persistence and variation, and student--job skill fit. Growth and stability will be labelled as proxy measures. Non-negative weights will sum to one; missing dimensions will be shown and handled through explicit coverage rules. Hybrid ranking will combine matching and quality scores without counting skill fit twice. Explanations will use the actual scores and matched skills. Keyword, skill-only, and hybrid approaches will be compared using HR@k, NDCG, MRR, and student trials.

  \item \textbf{Front-end Development}\par
  React and TypeScript will provide profile entry, region and qualification filters, weight controls, ranked results, score breakdowns, skill-gap views, occupation relationships, and feedback. The interface will call FastAPI through documented JSON contracts and show data dates, coverage, and empty-result guidance. Preference changes will be visible and reversible.
\end{enumerate}

\noindent\textbf{Engineering and delivery:}
GitHub reviews and automated checks will cover cleaning, score calculations, ranking, and API integration. Docker will package the application. Versioned data, dictionaries, scoring configurations, and evaluation splits will make recommendations reproducible. Our seven members will work across data, algorithms, and application delivery, with paired review connecting these workstreams.

\medskip
\noindent\textbf{Core data flow:}
Authorised recruitment data $\rightarrow$ cleaning and skill extraction $\rightarrow$ Parquet/DuckDB features $\rightarrow$ trends and occupation graph $\rightarrow$ profile-based matching and quality scoring $\rightarrow$ ranked explanations $\rightarrow$ user feedback.

\newpage
\section{Timeline \& Workflow and Work Distribution}
\subsection{Timeline \& Workflow}
Ziqian XIONG will lead weekly progress and integration reviews during concentrated development and use brief asynchronous check-ins in the low-workload period. Paired discussions of methods and code, and early demonstrations, will support supervisor feedback. The leader will help members break down unfamiliar tasks and resolve dependencies. Each member will prepare questions, report blockers early, and record decisions and follow-up actions.

The project runs from 21 October 2026 to final submission on 1 April 2027, with work organised by date ranges, workload, and demonstrable outputs. We use 16 December 2026--5 January 2027 to review documents and finish earlier tasks as needed. As January and February may overlap the winter vacation, 6 January--28 February is a low-workload period: at most one small asynchronous task per member per week, adjusted to availability, with no critical integration, required training runs, or participant trials. University holiday dates remain subject to formal notices. We bring the first demonstrable workflow and pilot experiments forward to 15 December, complete integration and evaluation in March, and reserve 31 March--1 April for submission checks.

Course milestones follow the supplied tentative key-dates table, subject to formal course notices.

\begin{figure}[!ht]
\centering
{\small
\setlength{\tabcolsep}{4pt}
\renewcommand{\arraystretch}{1.3}
\resizebox{\textwidth}{!}{%
\begin{tabular}{|l|*{7}{c|}}
\hline
\textbf{Phase} & \multicolumn{3}{c|}{2026} & \multicolumn{4}{c|}{2027} \\
\hline
Planning window & Oct & Nov & Dec & Jan & Feb & Mar & Apr \\
\hline
1. Requirements / data audit & \cellcolor{cyan!45}21+ & \cellcolor{cyan!45}1--3 & \rule{0pt}{1.6ex} & \rule{0pt}{1.6ex} & \rule{0pt}{1.6ex} & \rule{0pt}{1.6ex} & \rule{0pt}{1.6ex} \\
2. Cleaning / matching baseline & \rule{0pt}{1.6ex} & \cellcolor{cyan!45}4--24 & \rule{0pt}{1.6ex} & \rule{0pt}{1.6ex} & \rule{0pt}{1.6ex} & \rule{0pt}{1.6ex} & \rule{0pt}{1.6ex} \\
3. First workflow / pilot comparisons & \rule{0pt}{1.6ex} & \cellcolor{cyan!45}25+ & \cellcolor{cyan!45}1--15 & \rule{0pt}{1.6ex} & \rule{0pt}{1.6ex} & \rule{0pt}{1.6ex} & \rule{0pt}{1.6ex} \\
Review documents / finish earlier tasks & \rule{0pt}{1.6ex} & \rule{0pt}{1.6ex} & \cellcolor{gray!25}16+ & \cellcolor{gray!25}1--5 & \rule{0pt}{1.6ex} & \rule{0pt}{1.6ex} & \rule{0pt}{1.6ex} \\
4. Light review / documentation & \rule{0pt}{1.6ex} & \rule{0pt}{1.6ex} & \rule{0pt}{1.6ex} & \cellcolor{yellow!35}6+ & \cellcolor{yellow!35}1+ & \rule{0pt}{1.6ex} & \rule{0pt}{1.6ex} \\
5. Web / feedback integration & \rule{0pt}{1.6ex} & \rule{0pt}{1.6ex} & \rule{0pt}{1.6ex} & \rule{0pt}{1.6ex} & \rule{0pt}{1.6ex} & \cellcolor{cyan!45}1--9 & \rule{0pt}{1.6ex} \\
6. Evaluation / student trials & \rule{0pt}{1.6ex} & \rule{0pt}{1.6ex} & \rule{0pt}{1.6ex} & \rule{0pt}{1.6ex} & \rule{0pt}{1.6ex} & \cellcolor{cyan!45}10--16 & \rule{0pt}{1.6ex} \\
7. Release / report / demo & \rule{0pt}{1.6ex} & \rule{0pt}{1.6ex} & \rule{0pt}{1.6ex} & \rule{0pt}{1.6ex} & \rule{0pt}{1.6ex} & \cellcolor{cyan!45}17--30 & \rule{0pt}{1.6ex} \\
Submission checks & \rule{0pt}{1.6ex} & \rule{0pt}{1.6ex} & \rule{0pt}{1.6ex} & \rule{0pt}{1.6ex} & \rule{0pt}{1.6ex} & \cellcolor{cyan!45}31 & \cellcolor{cyan!45}1 \\
\hline
\end{tabular}}
}
\caption{Project timetable (dates within each month; \texttt{+} means through month end; grey: document review / earlier tasks; amber: low workload)}
\label{fig:timetable}
\end{figure}

\begin{itemize}
  \Needspace{3\baselineskip}
  \item \textbf{21 October -- 3 November 2026}
    \begin{itemize}
      \item Agree student journeys, acceptance criteria, authorised recruitment-data access, and the initial Zhejiang subset. Audit sample fields, dates, missing values, and duplicate records; draft the data dictionary and evaluation split.
      \item Measure sample-processing memory and runtime, and prepare storage/compute requests from 21 October. Submit ethics forms by 28 October, covering student profiles, feedback, consent, and planned user trials; begin participant work only after the required approval.
    \end{itemize}

  \Needspace{3\baselineskip}
  \item \textbf{4--24 November 2026}
    \begin{itemize}
      \item Publish the team project website by 5 November, with the project scope, team roles, and progress plan. Design Parquet/DuckDB data storage, profile tables, and FastAPI contracts; use synthetic profiles for interface development.
      \item Implement batch cleaning, salary-period normalisation, cross-platform deduplication, skill dictionaries, and keyword matching. Deliver a queryable authorised subset, quality statistics, and a reproducible recommendation baseline by 24 November. Prepare interface skeletons and reusable evaluation fixtures in parallel with this baseline.
    \end{itemize}

  \Needspace{3\baselineskip}
  \item \textbf{25 November -- 15 December 2026}
    \begin{itemize}
      \item Implement salary, growth, stability, and skill-fit scoring, chronological trend baselines, and a queryable occupation--skill graph. By 15 December, demonstrate profile entry, qualification/location filters, adjustable weights, ranked jobs, matched skills, skill gaps, and explanations on an authorised subset.
      \item In parallel, prepare React/TypeScript and FastAPI interface skeletons, a minimal explicit-feedback example, and pilot keyword, skill-only, and semantic comparisons with reusable evaluation scripts. Submit the interim report by 10 December, using evidence available then; archive code/data versions and integration issues by 15 December.
    \end{itemize}

  \Needspace{3\baselineskip}
  \item \textbf{16 December 2026 -- 5 January 2027}
    \begin{itemize}
      \item Use this time to review documents, collect feedback, and finish earlier tasks as needed; no critical work or participant trials are scheduled. Low-workload activity continues from 6 January.
    \end{itemize}

  \Needspace{3\baselineskip}
  \item \textbf{6 January -- 28 February 2027}
    \begin{itemize}
      \item Review existing skill dictionaries, data-quality notes, score explanations, and saved pilot results in small independent tasks; draft the API and user guides. These tasks must not block the March workflow.
      \item Use asynchronous check-ins to maintain the issue list and prepare test labels, trial instructions, and March priorities from existing materials. Make no full-application delivery commitment in January or February; resume concentrated integration on 1 March.
    \end{itemize}

  \Needspace{3\baselineskip}
  \item \textbf{1--9 March 2027}
    \begin{itemize}
      \item Complete the existing profile, ranking, score-breakdown, and graph views, keeping data dates and coverage visible. Connect saved preferences, favourites, and explicit feedback to reversible, user-approved weight updates; finalise versioned API and scoring records.
      \item By 9 March, demonstrate entry, recommendation, explanation, feedback, and refresh. Admit semantic matching only if the pilot supports it; keep a working baseline and add no complex model scope.
    \end{itemize}

  \Needspace{3\baselineskip}
  \item \textbf{10--16 March 2027}
    \begin{itemize}
      \item Use the prepared chronological splits, labels, and scripts to compare keyword, skill-only, and hybrid ranking with HR@k, NDCG, and MRR. Review skill extraction, trend errors, graph usefulness, and the contribution of each quality dimension.
      \item Run approved student trials and latency/resource checks in parallel; prioritise explanation clarity, empty results, and abnormal weights. Complete the evaluation report and essential fixes by 16 March.
    \end{itemize}

  \Needspace{3\baselineskip}
  \item \textbf{17--30 March 2027}
    \begin{itemize}
      \item Freeze features and models on 17 March, regress core journeys, and verify Docker deployment in a fresh environment by 23 March. Reproduce representative evaluation results and complete authorised-data access, configuration, limitation, and user instructions.
      \item Use 24--30 March to rehearse weight changes, skill gaps, and career-path demonstrations; complete the software release, source repository, final report, user guide, and demonstration materials by 30 March.
    \end{itemize}

  \Needspace{3\baselineskip}
  \item \textbf{31 March -- 1 April 2027}
    \begin{itemize}
      \item Verify release/report/documentation consistency, filenames, access, and receipts; submit the team final report and software by 1 April. Prepare demonstration, presentation, and promotional assets for later course submissions.
    \end{itemize}

\end{itemize}

\clearpage

\subsection{Work Distribution}
We will divide the work according to our relevant experience and interests, help one another learn unfamiliar methods, and adjust tasks through progress discussions and supervisor feedback.
\begin{itemize}
  \item \textbf{Ziqian XIONG:} Coordinate the team and supervisor discussions; guide system design, skill matching, ranking experiments, and model integration, drawing on his research and software projects.
  \item \textbf{Yun Xin NG:} Prepare and clean recruitment data, normalise salary fields, and develop skill-trend forecasts, drawing on her data-engineering and machine-learning experience.
  \item \textbf{Jinyao LU:} Organise student requirements and trials, check whether recommendation explanations are understandable, and prepare user documentation and demonstrations.
  \item \textbf{Oswald Yiwen CHONG:} Develop backend services and connect application modules; support GitHub collaboration, automated checks, and deployment using his application-development and CI/CD experience.
  \item \textbf{Ziqian ZHANG:} Prepare checked skill and occupation samples, organise annotations, and record data and recommendation errors, building on her research data-preparation experience.
  \item \textbf{Yue WU:} Develop skill-normalisation rules, four-dimensional scoring, and occupation relationships; compare methods and document the reasoning behind scores.
  \item \textbf{Xinyun HU:} Build student-profile, preference, and recommendation pages; connect them to backend services and check filtering, weight changes, and feedback interactions.
\end{itemize}

Yun Xin NG and Ziqian ZHANG will check data samples together; Yue WU and Ziqian XIONG will discuss scoring and ranking results. The application and user-trial work will connect these results to the student experience.

\clearpage
\section{Conclusion}
SmartCareer needs a team that can connect real recruitment data, transparent evaluation, and a usable student application. Our complementary preparation in data extraction, model experiments, programming, databases, delivery automation, and user coordination provides a credible basis for that work. We will demonstrate a reproducible baseline first, then use ranking metrics and student trials to justify improvements. This makes our bid a practical route from an available dataset to an explainable employment-support tool.

Building on our academic preparation, we aim to distinguish our team through stronger leadership and peer guidance, proactive learning, and consistent follow-through throughout the project.

\end{document}

